\scanpage{151}
% archive sheet number: 150
Supposons que l'on ait un système projectif $\unsure{\mathscr S|\mathcal S}=(S_i)_{i\in I}$, avec $i_0\in I$ tel que $i\geq i_0$ implique $\varphi_{i_0i}:S_i\unsure{\text{---}|\longrightarrow}S_{i_0}$ est affine. Soit $S=\varprojlim S_i$. Soit $G$ sur $S$, [d'un des trois types envisagés]. Dans beaucoup de situations, on dispose de mieux que d'un groupe sur $S$, mais on dispose \unsure{d'un} objet de
\[
\unsure{\mathcal V|\mathscr V}_\ell(\mathscr S)=\underset{\longrightarrow}{\operatorname{Pseudolim}}\mathcal V_\ell(S_i)
\]
[resp. \unsure{en|avec} $\mathcal V_\ell$ remplacé par $\unsure{\mathcal M|\mathscr M}_\ell$, $\unsure{\mathcal I|\mathrm I}_\ell$].

Exemple : Si par exemple, on part d'un \deleted{\unsure{prém}} $S$-schéma abélien \inserted{$A$}, d'où pour $i$ grand un $S_i$-schéma abélien $A_i$, d'où des $T_\ell(A_i)\in\mathcal M_\ell(S_i)$, qui définissent un élément de $\deleted{\unsure{\mathcal V_\ell}}\mathcal M_\ell(\mathscr S)$ [$\ell$ inversible \unsure{en|sur} $S$].
% slanted
\margin{Autre exemple : $X$ schéma \unsure{lisse} et propre sur $S$, $\ell$ premier aux car. rés. de $S$, on prend les faisceaux de Weil $\varprojlim R^if_*(\mathbb Z/\ell^{\unsure{n|u}}\mathbb Z)=R^if_*(\unsure{\mathbb Z_\ell|\mathbb Q_\ell})$ \unsure{mod} torsion\unsure{......}}

Supposons pour simplifier les $S_i$ connexes, donc $S$ connexe, soit $\xi$ un pt
